\clearpage
\section{Исследовательская часть}

Графовые модели построены по псевдокодам листингов~\ref{lst:iter}
и~\ref{lst:rec}. Вершина графа управления и информационного графа обозначена
номером строки оператора. Вершина истории содержит номер строки оператора
и пояснение срабатывания. Истории построены для входных данных $k=1$, $s=1$,
при которых ответ равен 1 (число 1).

\subsection{Графовые модели нерекурсивного алгоритма}

Граф управления нерекурсивного алгоритма показан на рисунке~\ref{fig:iter-control}.
Из начала цикла управление передаётся либо в тело цикла, либо
к оператору, следующему за циклом.

\begin{figure}[H]
    \centering
    \graphmodel{iter-control}
    \caption{Граф управления нерекурсивного алгоритма}
    \label{fig:iter-control}
\end{figure}

Информационный граф нерекурсивного алгоритма показан на
рисунке~\ref{fig:iter-info}. Петли у вершин 4--6 соответствуют вычислению
следующего значения переменной цикла по предыдущему, петля у вершины 7~---
накоплению суммы в ячейке таблицы.

\begin{figure}[H]
    \centering
    \graphmodel{iter-info}
    \caption{Информационный граф нерекурсивного алгоритма}
    \label{fig:iter-info}
\end{figure}

Операционная история нерекурсивного алгоритма показана на
рисунке~\ref{fig:iter-op-history}. Для начала цикла указано значение
переменной цикла; срабатывание, при котором условие продолжения ложно, является
последним срабатыванием начала цикла. Для оператора 7 указана изменяемая
ячейка таблицы.

\begin{figure}[H]
    \centering
    \graphmodel{iter-op-history}
    \caption{Операционная история нерекурсивного алгоритма}
    \label{fig:iter-op-history}
\end{figure}

Информационная история нерекурсивного алгоритма показана на
рисунке~\ref{fig:iter-info-history}. Срабатывания операторов 2--5 в начале
выполнения не используют результатов других срабатываний.

\begin{figure}[H]
    \centering
    \graphmodel{iter-info-history}
    \caption{Информационная история нерекурсивного алгоритма}
    \label{fig:iter-info-history}
\end{figure}

\subsection{Графовые модели рекурсивного алгоритма}

Граф управления рекурсивного алгоритма показан на рисунке~\ref{fig:rec-control}.
Сплошные дуги от операторов 10, 16 и 17 соответствуют вызовам функций,
штриховые дуги от операторов 11 и 19~--- возврату управления из функции
к оператору вызывающей функции, который выполняется после вызова.

\begin{figure}[H]
    \centering
    \graphmodel{rec-control}
    \caption{Граф управления рекурсивного алгоритма}
    \label{fig:rec-control}
\end{figure}

Информационный граф рекурсивного алгоритма показан на
рисунке~\ref{fig:rec-info}. Дуги от операторов вызова 10, 16 и 17 соответствуют
передаче аргументов, дуги от операторов 11 и 19~--- передаче результата
функции. Точками отмечено разветвление общей линии дуг от оператора 16.

\begin{figure}[H]
    \centering
    \graphmodel{rec-info}
    \caption{Информационный граф рекурсивного алгоритма}
    \label{fig:rec-info}
\end{figure}

Операционная история рекурсивного алгоритма показана на
рисунке~\ref{fig:rec-op-history}. Для каждого срабатывания указан вызов
функции, в котором оно выполнено.

\begin{figure}[H]
    \centering
    \graphmodel{rec-op-history}
    \caption{Операционная история рекурсивного алгоритма}
    \label{fig:rec-op-history}
\end{figure}

Информационная история рекурсивного алгоритма показана на
рисунке~\ref{fig:rec-info-history}. Срабатывания операторов 2, 7 и 10 в вызове
$F(1,1)$ используют только входные данные, а срабатывания операторов 4, 6
и 14 присваивают постоянные значения, поэтому входящих дуг у них нет.

\begin{figure}[H]
    \centering
    \graphmodel{rec-info-history}
    \caption{Информационная история рекурсивного алгоритма}
    \label{fig:rec-info-history}
\end{figure}

\subsection{Сравнение алгоритмов по графовым моделям}

Число вершин операционной истории равно числу срабатываний операторов,
число ярусов информационной истории~--- наименьшему числу последовательных
шагов при неограниченном числе исполнителей. Эти характеристики для
нескольких входных данных приведены в таблице~\ref{tbl:stats}; они получены
построением историй по псевдокодам листингов~\ref{lst:iter} и~\ref{lst:rec}.

\begin{table}[H]
\caption{Характеристики историй алгоритмов}
\label{tbl:stats}
\setlength{\extrarowheight}{2pt}
\begin{tabularx}{\linewidth}{|c|c|>{\centering\arraybackslash}X|>{\centering\arraybackslash}X|>{\centering\arraybackslash}X|>{\centering\arraybackslash}X|}
\hline
\multirow{2}{*}{$k$} & \multirow{2}{*}{$s$} &
\multicolumn{2}{c|}{Нерекурсивный алгоритм} &
\multicolumn{2}{c|}{Рекурсивный алгоритм} \\
\cline{3-6}
& & срабатываний & ярусов & срабатываний & ярусов \\
\hhline{|=:=:=:=:=:=|}
\inputrows graphs/stats.tex \end{tabularx}
\end{table}

Число срабатываний нерекурсивного алгоритма пропорционально числу
слагаемых в формуле~\eqref{eq:f}, то есть не превосходит величины порядка
$10ks$. В рекурсивном алгоритме значения $F(i,j)$ с одинаковыми аргументами
вычисляются повторно, поэтому число срабатываний растёт значительно быстрее:
при $k=4$, $s=18$ оно почти в 45 раз больше, чем у нерекурсивного алгоритма.
Число ярусов информационной истории у алгоритмов близко, а при б\'{о}льших
входных данных у рекурсивного алгоритма оно меньше. У нерекурсивного
алгоритма его увеличивают цепочки срабатываний начала цикла, каждое
из которых использует предыдущее значение переменной цикла.

\subsection{Фрагменты, допускающие параллельное выполнение}

В нерекурсивном алгоритме параллельно могут выполняться:
\begin{itemize}
    \item заполнение нулями разных ячеек таблицы в операторе 2;
    \item вычисление ячеек одной строки таблицы при фиксированном $i$: значение
    $F(i,j)$ использует только строку $i-1$ и записывается в собственную ячейку,
    поэтому итерации цикла по $j$ информационно независимы;
    \item сложение до десяти слагаемых суммы в формуле~\eqref{eq:f}, если
    накапливать частичные суммы отдельно и складывать их попарно; при таком
    малом числе слагаемых выигрыш не покрывает затрат на организацию
    параллельного выполнения.
\end{itemize}

В нерекурсивном алгоритме не могут выполняться параллельно:
\begin{itemize}
    \item вычисление разных строк таблицы: строка $i$ использует значения
    строки $i-1$, что видно по дугам информационной истории от срабатываний
    операторов 2 и 3, записавших строку 0, к срабатываниям оператора 7
    для ячейки $F[1][1]$;
    \item прибавление слагаемых к одной ячейке в операторе 7 без отдельных
    частичных сумм, так как каждое прибавление использует результат
    предыдущего;
    \item оператор 3 выполняется после заполнения ячейки $F[0][0]$ нулями,
    иначе записанное им значение будет потеряно;
    \item оператор 8 выполняется после вычисления ячейки $F[k][s]$.
\end{itemize}

Цепочки срабатываний начала цикла не ограничивают параллельное
выполнение: в параллельной реализации индексы ячеек могут быть заданы
каждому исполнителю непосредственно.

В рекурсивном алгоритме параллельно могут выполняться вызовы 16
и 17 одного вызова $S$: оба используют только аргументы этого вызова
и не связаны информационным отношением, что видно по общему предшественнику
и одинаковому ярусу их срабатываний в информационной истории. Тем самым
независимыми являются все ветви дерева рекурсивных вызовов.

В рекурсивном алгоритме не могут выполняться параллельно:
\begin{itemize}
    \item оператор 18 и вызовы 16 и 17, поскольку сумма требует
    результатов обоих вызовов;
    \item оператор 11 и вызов 10, поскольку оператор 11 возвращает результат вызова;
    \item проверки 2, 3, 7 и 13 и следующие за ними операторы, так как
    от результата проверки зависит выполнение остальных операторов вызова.
\end{itemize}

Каждый вызов выполняет лишь несколько операторов, поэтому при параллельной
реализации создание исполнителя для каждого вызова требует больше времени,
чем сам вызов; параллельно целесообразно выполнять только верхние уровни
дерева вызовов.
